% ------------------------------------------------------------------------------
% Slide 3: Evidencia en Nuestras Mediciones: Los Polos que Funcionan Sólido - U y A
% ------------------------------------------------------------------------------
\begin{frame}{Evidencia en Nuestras Mediciones: Los Polos que Funcionan Sólido - U y A}
    \begin{columns}[c]
        \begin{column}{0.51\textwidth}
            \scriptsize
            \begin{block}{Polo Labial - Vocal /u/ y /o/}
                \begin{itemize}\setlength{\itemsep}{1pt}\setlength{\parskip}{0pt}
                    \item El músculo de los labios (\textit{orbicularis}) hace trompita hacia adelante y alarga el tubo resonador.
                    \item Causa física directa: señal limpia, ruido bajo y respuesta predecible en todos los sujetos.
                \end{itemize}
            \end{block}
            \vspace{0.03cm}
            \begin{block}{Polo de Apertura Mandibular - Vocal /a/}
                \begin{itemize}\setlength{\itemsep}{1pt}\setlength{\parskip}{0pt}
                    \item Para subir $F_1$ necesitamos abrir la boca bajando la mandíbula.
                    \item Registro sólido: activación clara debajo del mentón (\textit{digástrico}) y firmeza en la mandíbula (\textit{masetero}).
                \end{itemize}
            \end{block}
            \vspace{0.03cm}
            \begin{alertblock}{El Desafío Pendiente}
                Con dos polos resueltos físicamente, queda resolver con la misma certeza el tercer vértice: \textbf{la vocal /i/}.
            \end{alertblock}
        \end{column}
        \begin{column}{0.47\textwidth}
            \centering
            \placeholderfigura{0.72\textheight}{Señales / Envolventes representativas mostrando la activación limpia de Orbicularis en U y Digástrico en A}
        \end{column}
    \end{columns}
\end{frame}
